\chapter{Introducci\'on}
\label{cap:introduccion}

\section{Contexto de la actividad}
\label{sec:contexto-actividad}

Esta actividad corresponde a la modalidad de trabajo de \textbf{GA (grupo
alterno)} de la Unidad IV de la asignatura \asignatura. A diferencia de las
entregas anteriores del curso, en esta modalidad el proyecto t\'ecnico
(PFC) no es desarrollado por el propio grupo, sino que se selecciona el PFC
de otro equipo ya terminado, y el grupo GA realiza sobre \'el un proceso de
\textbf{revisi\'on cruzada}: an\'alisis de arquitectura, evaluaci\'on de la
API REST y su documentaci\'on, investigaci\'on de tendencias relacionadas, y
retroalimentaci\'on cr\'itica independiente por parte de cada integrante
nuevo del grupo.

\section{Distinci\'on entre el equipo original del PFC y el \equipoga}
\label{sec:distincion-equipos}

Es importante dejar expl\'icita, desde el inicio de este informe, una
distinci\'on que se mantiene consistentemente en todo el documento:

\begin{itemize}[nosep]
  \item \textbf{Equipo original del PFC BIOPET (equipo BMT):} Mariscal
    Cabrera Jaime Josu\'e, Beltr\'an Montiel Fred Adri\'an y Taipe Mora
    Zaida Melissa. Este equipo dise\~n\'o, implement\'o y document\'o
    t\'ecnicamente BIOPET a lo largo de las entregas previas del curso
    (Entrega 1A, Entrega 1B, Tercera Entrega). El c\'odigo fuente, las
    pruebas automatizadas, la arquitectura y la evidencia t\'ecnica
    descrita en los cap\'itulos~\ref{cap:presentacion-pfc}
    y~\ref{cap:mvc-api-rest} de este informe son obra de ese equipo.
  \item \textbf{\equipoga\ (grupo de la actividad GA):} Mariscal Cabrera
    Jaime Josu\'e, Carvajal Loor Johan Stalin y Fajardo Montes Michael
    Xavier. Este es el grupo responsable de \textbf{este informe} y de la
    actividad de Unidad IV: selecci\'on del PFC, presentaci\'on,
    investigaci\'on comparativa, revisi\'on cr\'itica y consolidaci\'on de
    hallazgos.
\end{itemize}

Mariscal Cabrera Jaime Josu\'e es la \'unica persona presente en ambos
equipos: particip\'o en el desarrollo original de BIOPET y tambi\'en
integra el \equipoga. Carvajal Loor Johan Stalin y Fajardo Montes Michael
Xavier \textbf{no participaron en el desarrollo del PFC}: su rol en esta
actividad es exclusivamente de revisi\'on cr\'itica independiente e
investigaci\'on, nunca de implementaci\'on del c\'odigo de BIOPET. Ning\'un
aporte de c\'odigo o funcionalidad del PFC debe atribuirse a Carvajal o a
Fajardo en este informe.

\section{Objetivo de este informe}
\label{sec:objetivo-informe}

El objetivo de este documento es consolidar, en un \'unico informe
acad\'emico, los siguientes componentes exigidos por la actividad de
Unidad IV:

\begin{enumerate}[nosep]
  \item Presentaci\'on del PFC seleccionado (BIOPET): prop\'osito,
    arquitectura y tecnolog\'ias (cap\'itulo~\ref{cap:presentacion-pfc}).
  \item An\'alisis del patr\'on MVC y de la API REST, incluyendo su
    documentaci\'on Swagger/OpenAPI y la evidencia de endpoints probados
    (cap\'itulo~\ref{cap:mvc-api-rest}).
  \item Investigaci\'on comparativa SOAP vs REST
    (cap\'itulo~\ref{cap:soap-vs-rest}).
  \item Investigaci\'on de tendencias web actuales: Jamstack, PWA e IA
    generativa (cap\'itulo~\ref{cap:tendencias-web}).
  \item Revisi\'on cruzada consolidada del PFC, a partir de las revisiones
    independientes de Carvajal y Fajardo
    (cap\'itulo~\ref{cap:revision-cruzada}).
  \item Resultados y evidencias t\'ecnicas verificables del estado actual
    de BIOPET (cap\'itulo~\ref{cap:resultados}).
  \item Conclusiones y trabajo futuro
    (cap\'itulos~\ref{cap:conclusiones} y~\ref{cap:trabajo-futuro}).
\end{enumerate}

\section{Estado de este documento}
\label{sec:estado-documento}

Esta versi\'on del informe integra ya contenido definitivo en sus nueve
cap\'itulos: la evidencia t\'ecnica verificable de BIOPET
(cap\'itulos~\ref{cap:presentacion-pfc} y~\ref{cap:mvc-api-rest}), la
investigaci\'on comparativa SOAP vs REST
(cap\'itulo~\ref{cap:soap-vs-rest}), la investigaci\'on de tendencias web
(cap\'itulo~\ref{cap:tendencias-web}), la revisi\'on cruzada consolidada de
Carvajal y Fajardo (cap\'itulo~\ref{cap:revision-cruzada}), los resultados y
evidencias emp\'iricas (cap\'itulo~\ref{cap:resultados}), y las
conclusiones y el trabajo futuro derivados de todo lo anterior
(cap\'itulos~\ref{cap:conclusiones} y~\ref{cap:trabajo-futuro}). Ninguno de
estos cap\'itulos conserva marcadores de contenido pendiente.

Esta versi\'on corresponde al documento final integrado: el PDF ya fue
generado y compilado con MiKTeX (\texttt{biblatex}/\texttt{biber}, seg\'un
la cadena de compilaci\'on descrita en \texttt{docs/u4/informe/README.md}),
y su resultado fue revisado y validado para la entrega de esta actividad.
